%第3章
\newpage
\part{結論}

\section{まとめ}
本研究では泥濘地でのスタックを防止するため, 各車輪の方向を進行方向にそろえる機構をもつ8輪運搬ロボットの提案と製作を行った.
また現場ヤードを想定した検証を行った.
    
検証の結果, 平地走行や段差踏破では十分な性能を持つことが検証できた.
よって本機体の有用性を証明する事ができたと言える.


\section{今後の展望}
今後は以下に示した土地盤で含水率を変えつつ走行実験を行い, 泥濘地において車輪を使用した機体の有用性の評価と, 方向可変な車輪のスタック防止に対する有用性の評価を行う必要がある.

\begin{itemize}
  \item 礫
  \item 砂
  \item シルト
  \item 粘土
\end{itemize}

\section{発展的な展望}
本機体の車輪形状は, 想定される様々な環境に合う設計にしているが, あくまでも創造しうる範囲での設計である.

本研究で現場ヤードでの走行性能を示す事はできたが, 今回の車輪形状が最適である証明はできていない.最適な形状を証明するには, いくつかの形状を製作し比較実験を行う必要がある.
また, 今回泥濘地での走行実験は行えていない為, 泥濘地において最適な車輪形状の検討も必要である.

さらに泥濘地での実験結果しだいでは, 車輪を使用した機体の1番の難点である, 接地面積の少なさを改善する必要が出てくる可能性もある.
よって, 接地面積増加を目標とした新機構の考案も視野に入れて今後の研究を進める必要がある.





